\section{Capability evaluation}
\label{sec:capability}

This section asks whether the synthesized API is useful for
downstream LLM agents (\S\ref{sec:cap-cross},
\S\ref{sec:cap-rlsteelman}), holds across LLM tiers
(\S\ref{sec:cap-portability}), and transfers beyond SO-101
(\S\ref{sec:cap-cross-robot}). Same-API comparisons hold the driver
fixed (agent style); synthesis itself is evaluated only in
\S\ref{sec:cap-portability}.

\subsection{SO-101 cross-skill: Reach-and-Pick}
\label{sec:cap-cross}

For this experiment the SO-101 driver is produced in spec mode
(\S\ref{sec:modes}), which exposes two interfaces: a low-level
joint interface the PPO and Diffusion Policy baselines read directly,
and the high-level Cartesian+gripper API the LLM uses. To make the RL comparison
like-for-like, we also train standard PPO on that high-level API. Its
action space is a Cartesian delta resolved by the driver's IK, plus a
gripper bit. We goal-condition it in three regimes (reach-only,
pick-only, both), $N{=}3$ seeds each (Appendix~\ref{app:robustness}). Reach evaluates 4 variants $\times$
10 episodes; Pick 10 episodes in which a banana must be lifted
$\geq$2\,cm from a randomized XY pose (same-API rows: $n{=}30$).
Table~\ref{tab:cross-skill} reports per-policy success rates.

\begin{table}[t]
\centering
\small
\setlength{\tabcolsep}{3pt}
\begin{tabular}{lrr}
\toprule
\textbf{Policy} & \textbf{Reach} & \textbf{Pick} \\
\midrule
\multicolumn{3}{l}{\emph{Low-level joint interface (no-abstraction ablation)}}\\
PPO\_reach (multi-task, 300k) & 100\% & 20\% \\
PPO\_pick (stoch.\ own-task)  &   0\% & 70\% \\
DP\_reach (600 demos)         &  50\% &  0\% \\
DP\_pick (200 demos)          &   0\% & 10\% \\
\midrule
\multicolumn{3}{l}{\emph{Same high-level API as the LLM ($N{=}3$ seeds, $n{=}30$)}}\\
RL\_reach (one skill)         & \textbf{100\%} & 26\% \\
RL\_pick (one skill)          &   0\% & \textbf{100\%} \\
RL\_multi (both skills)       & \textbf{100\%} & \textbf{100\%} \\
CaP (zero-shot)               & \textbf{100\%} & \textbf{100\%} \\
\textbf{Auto-Adapter} (zero-shot) & \textbf{100\%} & \textbf{100\%} \\
\bottomrule
\end{tabular}
\caption{SO-101 cross-skill (Reach, Pick). \emph{Top:} learned policies on
the low-level joint interface (no-abstraction ablation). \emph{Bottom:}
PPO on the \emph{same} high-level API the LLM uses, plus the two LLM methods
(RL rows: $N{=}3$ seeds; CaP/Auto-Adapter zero-shot). PPO\_pick is the charitable
\emph{stochastic} own-task rate (deterministic 0\%, an SB3 artifact).
Top rows: \texttt{exp3\_cross\_skill}; same-API rows: \texttt{exp\_same\_api\_rl}
(canonical results file).}
\label{tab:cross-skill}
\end{table}

Both LLM methods reach 100\% on both skills zero-shot, and the
same-API learned policies isolate why. A single-skill PPO masters its
trained skill (100\%) but transfers poorly: run on Pick, the
reach-trained policy succeeds on only $26\pm13\%$ (3-seed mean)
versus 100\% for the pick-trained, and RL\_pick scores 0\% on
Reach. A multi-task PPO reaches 100\%/100\%, so the limit is the
training regime, not the API. RL matches the LLM only after training
on every target skill; the LLM covers new skills at zero added
training cost. The joint-interface policies (top) score lower still,
showing how much the synthesized API's abstraction contributes.

\subsection{Multi-task RL ablation}
\label{sec:cap-rlsteelman}

RL is not simply weak given the right data. Trained on 6 reach
directions and tested on 6 in-distribution plus 4 held-out variants,
multi-task PPO reaches 100\% / 75\% and our LLM agent 100\% / 75\%;
both miss the same 10\,cm out-of-workspace variant. Multi-task DP
reaches 60\% / 5\% (Appendix~\ref{app:multitask}). The cross-skill gap
(\S\ref{sec:cap-cross}) is thus not that PPO is weak on its own task,
but that it needs a new environment, reward, and training run per
skill family.

\subsection{Using vs.\ synthesizing the tool layer}
\label{sec:cap-portability}

We separate two questions:
\emph{using} a \emph{fixed} reference driver (synthesized once,
frozen for all models) vs.\
\emph{synthesizing} one's own from the MJCF. The comparison
spans seven LLMs across five vendors on the 13-task SO-101 suite
($N{=}5$, trial-level; adapters in Appendix~\ref{app:models}).
Table~\ref{tab:portability} reports both.

% Headline ablation: USING a synthesized tool layer vs SYNTHESIZING it.
% Numbers are trial-level physics success on the SO-101 13-task suite (N=5),
% corrected (numpy-bool serialization bug fixed) with the feasible
% waypoint-verified draw_letter_L. Bold = best of the three "uses" methods.
\begin{table}[t]
\centering
\small
\setlength{\tabcolsep}{2.5pt}
\begin{tabular}{llrrrcr}
\toprule
 & & \multicolumn{3}{c}{\textbf{Uses API}} & \multicolumn{2}{c}{\textbf{Writes API}} \\
\cmidrule(lr){3-5}\cmidrule(lr){6-7}
\textbf{Model} & \textbf{Vendor} & \textbf{A-A} & \textbf{CaP} & \textbf{CaP\textsubscript{fs}} & \textbf{Syn} & \textbf{Self} \\
\midrule
Sonnet~4.6    & Anthropic & \textbf{69} & 57          & 54 & $\times^{\,\mathrm{V}}$ & -- \\
Opus~4.8      & Anthropic & \textbf{63} & 51          & 42 & \checkmark & 46 \\
Haiku~4.5     & Anthropic & \textbf{62} & 54          & 62 & \checkmark & 46 \\
Nova~Pro      & Amazon    & \textbf{60} & 57          & 46 & $\times^{\,\mathrm{S}}$ & -- \\
DeepSeek~V3.2 & DeepSeek  & 57          & \phantom{0}9 & \textbf{60} & \checkmark & 28 \\
Ministral~8B  & Mistral   & \textbf{54} & 23          & 43 & $\times^{\,\mathrm{S}}$ & -- \\
Qwen3~32B     & Alibaba   & 49          & \textbf{72} & 57 & $\times^{\,\mathrm{V}}$ & -- \\
\bottomrule
\end{tabular}
\caption{\textbf{Using vs.\ synthesizing the tool layer} (SO-101,
13 tasks, pass \%, $N{=}5$). \textbf{Uses API} (fixed driver):
A-A = Auto-Adapter, CaP = zero-shot Code-as-Policies,
CaP\textsubscript{fs} = CaP with few-shot exemplars; bold = best of
the three. \textbf{Writes API}: Syn = passes validation
($\times^{\mathrm{S}}$/$\times^{\mathrm{V}}$ = study/validate
failure); Self = end-to-end on the model's own driver.}
\label{tab:portability}
\end{table}


\noindent\textbf{Using.} On the \emph{same} fixed driver,
Auto-Adapter tops zero-shot CaP on six of seven models (49--69\%
vs.\ 9--72\%); the gap reflects the multi-turn planner, not the
synthesis. Significance is at the task
level, not per model: a Wilcoxon test over 91 model-task cells
gives $p{<}0.01$; per-model tests are underpowered
(Appendix~\ref{app:robustness}). Hand-written few-shot
exemplars (CaP\textsubscript{fs}, extra human effort per API) only
rescue weak models that cannot use the API zero-shot (DeepSeek
$9{\to}60$) while \emph{hurting} strong ones (Sonnet~4.6 $57{\to}54$);
against CaP\textsubscript{fs} Auto-Adapter still wins four of
seven, trailing only on DeepSeek and Qwen.

\noindent\textbf{Synthesizing.} Using a tool layer does not imply
writing one. All seven operate the API ($\geq$49\%), but only three
synthesize a driver that passes validation (build, IK, gripper,
perception, grasp). Opus, Haiku, and DeepSeek pass; Sonnet and Qwen fail
validation (grasp, IK); Nova and Ministral fail earlier ($0/3$).
Qwen (best CaP user, 72\%) and Sonnet (best API user, 69\%)
both fail to ship a driver. Operating their self-written drivers
end-to-end (``Self''), the three reach 28--46\% --- below their own
scores on the fixed reference (57--63\%), but with \emph{zero}
human driver code.
Figure~\ref{fig:cost-quality} plots the cost--quality plane
(debugging-loop traces in Figure~\ref{fig:synth-traj}).

\emph{Practitioner takeaway: capability is needed only at
onboarding: synthesize once with a frontier model, then serve with
a cheap one.}

\subsection{Cross-robot scaling}
\label{sec:cap-cross-robot}

\noindent\textbf{Simple suite on Piper and UR5e.}
On five short arm tasks ($\times$3 trials; three physics-graded,
two behavioral, \S\ref{sec:protocol}), Auto-Adapter reaches 80\% and
CaP 60\% on each robot. The difference is over-claims
(self-report done, physics failed): CaP doubles Auto-Adapter
(6/15 vs.\ 3/15). All three extra over-claims fall on
\texttt{unreachable\_detect}, where CaP reports completion despite an
\texttt{IKUnreachableError} that Auto-Adapter catches
(filmstrips in Appendix~\ref{app:filmstrips}).

\noindent\textbf{End-to-end Pick on Piper.}
Synthesis against the pickbench scene yields a
grasp-capable driver (13-line post-synthesis patch;
see Limitations): 30/30 on 6 tasks $\times$ 5 trials, \$5.02 /
18.9\,min end-to-end.

\noindent\textbf{Quadruped locomotion.}
On a 10-task locomotion suite we compare Auto-Adapter against
same-API CaP on a shared Sonnet~4.5 driver with physics grading.
The outcome is embodiment-dependent: ANYmal-C favors Auto-Adapter
(68\% vs.\ 28\%), while Go2 reverses it (48\% vs.\ 54\%). CaP's
open-loop gait clears only one short ANYmal-C walk; on Go2's coarser
$\approx$0.15\,m step CaP's tuned walk transfers while Auto-Adapter
overshoots. The split holds under a Sonnet~4.6 agent
(Appendix~\ref{app:robustness}).
